\begin{mistake}{2026-06-29}{04}

\begin{problemcontent}
计算不定积分 \(\int \frac{\arcsin x}{x^2} \mathrm{d}x\)。
\end{problemcontent}

\begin{solutioncontent}
令 \(u = \arcsin x\)，\(\mathrm{d}v = \frac{1}{x^2} \mathrm{d}x\)，则 \(\mathrm{d}u = \frac{1}{\sqrt{1-x^2}} \mathrm{d}x\)，\(v = -\frac{1}{x}\)。
应用分部积分法：
\[
\int \frac{\arcsin x}{x^2} \mathrm{d}x = -\frac{\arcsin x}{x} + \int \frac{1}{x\sqrt{1-x^2}} \mathrm{d}x
\]
对于剩余积分，采用分子分母同乘 \(x\) 的技巧：
\[
\int \frac{x}{x^2\sqrt{1-x^2}} \mathrm{d}x
\]
令 \(t = \sqrt{1-x^2}\)，则 \(t^2 = 1 - x^2\)，\(x \mathrm{d}x = -t \mathrm{d}t\)。
\[
\int \frac{1}{x\sqrt{1-x^2}} \mathrm{d}x = \int \frac{-t}{(1-t^2)t} \mathrm{d}t = \int \frac{1}{t^2-1} \mathrm{d}t = \frac{1}{2}\ln\left|\frac{t-1}{t+1}\right| + C
\]
代回 \(t = \sqrt{1-x^2}\) 并有理化真数化简得：
\[
\frac{1}{2}\ln\left|\frac{\sqrt{1-x^2}-1}{\sqrt{1-x^2}+1}\right| = \ln\left|\frac{1-\sqrt{1-x^2}}{x}\right|
\]
故最终结果为：
\[
\int \frac{\arcsin x}{x^2} \mathrm{d}x = -\frac{\arcsin x}{x} + \ln\left|\frac{1-\sqrt{1-x^2}}{x}\right| + C
\]
\end{solutioncontent}

\mistakeknowledge{
\begin{itemize}
  \item \textbf{核心公式}：分部积分 \(\int u \mathrm{d}v = uv - \int v \mathrm{d}u\)；\(\int \frac{1}{t^2-a^2} \mathrm{d}t = \frac{1}{2a}\ln\left|\frac{t-a}{t+a}\right| + C\)
  \item \textbf{易错点}：求解含 \(\frac{1}{x\sqrt{a^2-x^2}}\) 的无理积分时，分子分母同乘 \(x\) 凑微分（或使用倒代换 \(x=1/t\)）是最佳策略，比三角换元更简便。
  \item \textbf{关联知识}：对数结果的标准形式化简，需熟练掌握分子分母同乘共轭根式的有理化技巧。
\end{itemize}}

\end{mistake}
